\section{\texorpdfstring{PPE 面係数 $a_f = 2/(\rho_L+\rho_R)$ の FVM 積分論理と周期境界対応}{PPE 面係数 a\_f = 2/(ρ\_L+ρ\_R) の FVM 積分論理と周期境界対応}}
\label{app:fvm_face_coeff}
% ═══════════════════════════════════════════════════════════════════════════════

PPE 演算子の係数 $a(x) = 1/\rho(x)$ について，
セル $P$（位置 $x_L$）と $E$（位置 $x_R$）の間の等価面係数を FVM 積分で求める．

フラックス $F = a_f(p_E - p_P)/\Delta x$ において，
連続体のフラックス $F = a(x)\,dp/dx$ を区間 $[x_L, x_R]$ で積分して面係数 $a_f$ を求める．
$a(x)$ が区間内で区分的定数（各セル内で $a_L = 1/\rho_L$，$a_R = 1/\rho_R$）のとき：
\[
  \int_{x_L}^{x_R} \frac{dx}{a(x)}
  = \frac{\Delta x/2}{a_L} + \frac{\Delta x/2}{a_R}
  \quad\Rightarrow\quad
  a_f = \frac{\Delta x}{\dfrac{\Delta x/2}{a_L}+\dfrac{\Delta x/2}{a_R}}
  = \frac{2a_L a_R}{a_L+a_R}
  = \frac{2}{\rho_L+\rho_R}
\]
これは直列抵抗の等価抵抗（$1/R_\mathrm{eq} = 1/R_1+1/R_2$）と同じ論理（調和平均）であり，
熱伝導率の界面補間で使われる標準公式と同一構造を持つ．\qed

\subsection{FVM PPE 行列の周期境界条件対応}
\label{sec:fvm_periodic}
%% 旧所在: §8.4（08f_ppe_bc.tex）→ 付録 E.4 へ移動（CHK-086）

壁面（ノイマン）BC に対する FVM 離散化（式~\eqref{eq:af_exact}）は，
周期境界を持つ計算領域ではそのまま適用できない．
本稿の実装では以下の3点を変更して周期 FVM 行列を組み立てる．

\paragraph{（1）折り返しフェイスの追加}
周期 BC では内部フェイスに加え，ノード $N_x-1$ とノード $0$ を結ぶ
\textbf{折り返しフェイス（wrap face）}を追加する：
\begin{equation}
  a_{N_x - 1/2} = \frac{2}{\rho_{N_x-1} + \rho_0}
  \label{eq:af_wrap}
\end{equation}

\paragraph{（2）境界半格子体積補正の省略}
周期 BC ではすべてのノードが同じ格子幅 $h$ の制御体積を持つため，
壁面 BC の $h/2$ 体積補正は行わない（$a_f / h^2$ を一様に適用）．

\paragraph{（3）ゴーストノードの同一性拘束}
インデックス $N$ ノード（物理的に $0$ の周期像）を
\textbf{同一性拘束行}で置き換える：
\begin{equation}
  p_{\mathrm{ghost}} - p_{\mathrm{src}} = 0
  \quad\Longleftrightarrow\quad
  A[\mathrm{ghost},\,\mathrm{ghost}] = 1,\;
  A[\mathrm{ghost},\,\mathrm{src}]   = -1,\;
  b[\mathrm{ghost}] = 0
  \label{eq:ghost_identity}
\end{equation}
コーナー節点は最初の軸の拘束のみを採用する（重複適用を避けるため）．

\begin{tcolorbox}[warnbox, title={周期 FVM 行列と前処理：ILU の充填不足に注意}]
  式~\eqref{eq:ghost_identity} の同一性拘束行と FVM 行が混在した疎行列に対して，
  ILU(0)（\texttt{fill\_factor=1}）を前処理に使用すると BiCGSTAB の収束が著しく遅延する．
  実装では \texttt{fill\_factor=10} で安定収束を確認した．
  直接法（\texttt{spsolve}）は収束診断に有効である．
\end{tcolorbox}

